\documentclass[UTF8,a4paper,12pt,fontset=none]{ctexart}

\usepackage{amsmath}
\usepackage{amssymb}
\usepackage{booktabs}
\usepackage{geometry}
\usepackage{xcolor}
\usepackage{listings}

\geometry{left=2.4cm,right=2.4cm,top=2.5cm,bottom=2.5cm}
\setCJKmainfont{PingFang SC}
\setCJKsansfont{Heiti SC}
\setCJKmonofont{PingFang SC}

\lstdefinestyle{hipstyle}{
  language=C++,
  basicstyle=\ttfamily\small,
  keywordstyle=\color{blue!70!black},
  commentstyle=\color{green!40!black},
  stringstyle=\color{red!60!black},
  numbers=left,
  numberstyle=\tiny\color{gray},
  stepnumber=1,
  numbersep=8pt,
  breaklines=true,
  frame=single,
  tabsize=4,
  showstringspaces=false
}
\lstset{style=hipstyle}

\title{AMD HIP 实验 Notebook 填空报告}
\author{}
\date{}

\begin{document}
\maketitle

\section{Lab 1：Vector Addition}

\subsection{线程索引计算}

一维网格中的全局线程编号为
\[
  \text{globalIdx}=\text{blockIdx.x}\times \text{blockDim.x}+\text{threadIdx.x}.
\]

\begin{center}
\begin{tabular}{cccc}
\toprule
blockIdx.x & blockDim.x & threadIdx.x & Global Index \\
\midrule
0 & 256 & 100 & 100 \\
2 & 128 & 50 & 306 \\
5 & 64 & 0 & 320 \\
\bottomrule
\end{tabular}
\end{center}

\subsection{Vector Add Kernel 填空}

\begin{lstlisting}
__global__ void vector_add(const float* A, const float* B, float* C, int N) {
    int idx = blockIdx.x * blockDim.x + threadIdx.x;

    if (idx < N) {
        C[idx] = A[idx] + B[idx];
    }
}
\end{lstlisting}

\subsection{启动配置计算}

当 $N=1000$ 时：

\begin{center}
\begin{tabular}{ccccc}
\toprule
Block Size & Grid Size & Total Threads & Idle Threads & Utilization \\
\midrule
64  & 16 & 1024 & 24  & 97.66\% \\
128 & 8  & 1024 & 24  & 97.66\% \\
256 & 4  & 1024 & 24  & 97.66\% \\
512 & 2  & 1024 & 24  & 97.66\% \\
\bottomrule
\end{tabular}
\end{center}

这四种 block size 在该例中的线程利用率相同。实际性能不一定完全由线程利用率决定，还会受到 wavefront 对齐、occupancy、寄存器使用、内存带宽和调度开销影响。Vector Add 的计算量很小，主要瓶颈是全局内存读写带宽，因此不同 block size 的性能差异通常不会像计算密集型 kernel 那样明显。

\subsection{Wavefront 分析}

默认 wave32 模式下，block size 为 100 时：
\[
  \left\lceil\frac{100}{32}\right\rceil=4
\]
因此每个 block 需要 4 个 wavefront，浪费 lane 数为
\[
  4\times 32-100=28.
\]

若使用 wave64：
\[
  \left\lceil\frac{100}{64}\right\rceil=2,
  \qquad 2\times 64-100=28.
\]
即需要 2 个 wavefront，同样浪费 28 个 lane。

当同一个 wavefront 内线程进入不同 if-else 分支时，会发生分支发散。硬件通常需要串行执行不同分支路径，并屏蔽不属于当前路径的 lane，因此有效并行度下降。

\subsection{Occupancy 分析}

occupancy 表示一个 CU 上活跃 wavefront 数与最大可支持 wavefront 数的比例。较高 occupancy 有助于隐藏内存延迟，但并不保证性能最优。若 kernel 受内存带宽、atomic contention 或指令依赖限制，继续增加 occupancy 可能无法带来明显收益。

\subsection{实测性能分析}

在 $N=1,048,576$、10 次运行下，实测结果如下：

\begin{center}
\begin{tabular}{ccccc}
\toprule
Block Size & Mean (ms) & Std (ms) & Min (ms) & Max (ms) \\
\midrule
16   & 0.0355 & 0.0009 & 0.0329 & 0.0359 \\
32   & 0.0201 & 0.0009 & 0.0187 & 0.0207 \\
64   & 0.0193 & 0.0015 & 0.0177 & 0.0223 \\
128  & 0.0202 & 0.0015 & 0.0176 & 0.0226 \\
256  & 0.0208 & 0.0018 & 0.0177 & 0.0252 \\
512  & 0.0207 & 0.0000 & 0.0206 & 0.0208 \\
1024 & 0.0212 & 0.0014 & 0.0207 & 0.0253 \\
\bottomrule
\end{tabular}
\end{center}

最低平均执行时间由 block size 64 取得，为 0.0193 ms。Vector Add 每个元素读取两个 float 并写入一个 float，数据量约为 $3N\times4=12,582,912$ bytes，因此 block size 64 的估算带宽为
\[
  \frac{12,582,912}{0.0193\times10^{-3}}\approx 651.96\ \text{GB/s}.
\]
由于带宽和执行时间成反比，本组数据中最高估算带宽同样对应 block size 64。除 block size 16 明显较慢外，32 到 1024 的误差范围有重叠，差异不宜解释为强统计显著。Vector Add 性能主要由内存带宽决定，因为每个元素只有一次加法，计算量远小于内存读写量。

\subsection{Sub-wavefront Block Size}

sub-wavefront 实测结果如下：

\begin{center}
\begin{tabular}{cccc}
\toprule
Block Size & Mean (ms) & Std (ms) & Sub-Wave? \\
\midrule
8   & 0.0551 & 0.0015 & YES \\
16  & 0.0350 & 0.0035 & YES \\
32  & 0.0211 & 0.0014 & no \\
64  & 0.0217 & 0.0049 & no \\
128 & 0.0206 & 0.0001 & no \\
256 & 0.0206 & 0.0001 & no \\
\bottomrule
\end{tabular}
\end{center}

block size 8 和 16 小于 wave32，会造成一个 wavefront 中只有部分 lane 有效工作。实测中 8 和 16 的平均时间分别为 0.0551 ms 和 0.0350 ms，明显慢于 wave-aligned 的 32、128、256，因此存在明显 sub-wavefront 性能损失。若换成 MI100 的 wave64，则所有小于 64 的 block size 都属于 sub-wavefront，例如 8、16、32。

\subsection{多元素每线程 Kernel}

\begin{lstlisting}
__global__ void vector_add_multi(const float* A, const float* B, float* C,
                                 int N, int elements_per_thread) {
    int tid = blockIdx.x * blockDim.x + threadIdx.x;
    int stride = blockDim.x * gridDim.x;

    for (int i = tid; i < N; i += stride) {
        C[i] = A[i] + B[i];
    }
}
\end{lstlisting}

相比一线程处理一个元素，grid-stride loop 能让固定数量的线程覆盖任意大的输入，减少启动过多 block 的需求，并在大规模数据上提高每个线程的工作量和调度效率。

\section{Lab 2：Matrix Transpose}

\subsection{二维索引}

二维索引计算为
\[
  idx=\text{blockIdx.x}\times \text{blockDim.x}+\text{threadIdx.x},
  \qquad
  idy=\text{blockIdx.y}\times \text{blockDim.y}+\text{threadIdx.y}.
\]

当 \texttt{blockDim=(32,32)} 时：

\begin{center}
\begin{tabular}{ccc}
\toprule
blockIdx & threadIdx & Global $(idx,idy)$ \\
\midrule
(0,0) & (5,10)  & (5,10) \\
(1,0) & (0,0)   & (32,0) \\
(2,3) & (15,20) & (79,116) \\
\bottomrule
\end{tabular}
\end{center}

对于 \texttt{rows=100, cols=200} 且 block size 为 $32\times32$ 的矩阵：
\[
  \text{blocks.x}=\left\lceil\frac{200}{32}\right\rceil=7,
  \qquad
  \text{blocks.y}=\left\lceil\frac{100}{32}\right\rceil=4.
\]
因此总 block 数为 $7\times4=28$。

\subsection{转置索引变换}

对于 $4\times6$ 矩阵：

\begin{center}
\begin{tabular}{ccc}
\toprule
Thread $(idx,idy)$ & Input Index & Output Index \\
\midrule
(0,0) & $0\times6+0=0$ & $0\times4+0=0$ \\
(1,0) & $0\times6+1=1$ & $1\times4+0=4$ \\
(0,1) & $1\times6+0=6$ & $0\times4+1=1$ \\
(3,2) & $2\times6+3=15$ & $3\times4+2=14$ \\
\bottomrule
\end{tabular}
\end{center}

写入索引使用 \texttt{idx * rows + idy}，是因为转置后输出矩阵尺寸为 \texttt{cols x rows}，输出矩阵每一行包含 \texttt{rows} 个元素，而不是原矩阵的 \texttt{cols} 个元素。

\subsection{Naive Transpose Kernel}

\begin{lstlisting}
__global__ void matrix_transpose_kernel(const float* input, float* output,
                                        int rows, int cols) {
    int idx = blockIdx.x * blockDim.x + threadIdx.x;
    int idy = blockIdx.y * blockDim.y + threadIdx.y;

    if (idx < cols && idy < rows) {
        output[idx * rows + idy] = input[idy * cols + idx];
    }
}
\end{lstlisting}

\subsection{Coalescing 分析}

读取 \texttt{input[idy * cols + idx]} 时，同一行内连续的 \texttt{threadIdx.x} 会访问连续地址，因此读取是 coalesced 的。

写入 \texttt{output[idx * rows + idy]} 时，连续 \texttt{threadIdx.x} 的写入地址相差 \texttt{rows} 个 float。对于 $1024\times1024$ 矩阵，stride 为 $1024$ 个 float，即 $4096$ bytes。

若一次内存 transaction 获取 128 bytes，也就是 32 个 float，则一个 wave32 做这种非连续写入时，理想情况下每个线程可能落在不同 transaction 中，因此约需要 32 次 transaction。

\subsection{Tiled Transpose 分析}

共享内存 tile 定义为 \texttt{tile[BLOCK\_SIZE][BLOCK\_SIZE + 1]}，其中 \texttt{+1} 用于打破列访问时的 bank conflict，使线程访问共享内存列时不会集中到同一个 bank。

写回时交换 \texttt{blockIdx.x} 和 \texttt{blockIdx.y}，是因为 tile 需要被转置到输出矩阵的对称位置，从而让读和写都尽量保持连续访问。

当 \texttt{BLOCK\_SIZE=32} 时，共享内存使用量为
\[
  32\times(32+1)\times4=4224\ \text{bytes}.
\]

\subsection{实测性能分析}

矩阵转置的实测 kernel 时间如下：

\begin{center}
\begin{tabular}{ccccc}
\toprule
Test & Dimensions & Elements & Mean (ms) & Std (ms) \\
\midrule
1 & $16\times16$     & 256       & 0.0078 & 0.0001 \\
2 & $128\times32$    & 4096      & 0.0090 & 0.0000 \\
3 & $1\times1024$    & 1024      & 0.0076 & 0.0010 \\
4 & $1001\times2001$ & 2,003,001 & 0.1117 & 0.0002 \\
\bottomrule
\end{tabular}
\end{center}

Test 4 的元素数量约为 Test 2 的 $2,003,001/4096\approx489$ 倍，但平均时间只从 0.0090 ms 增加到 0.1117 ms，约为 $12.4\times$。时间没有按元素数线性增长，原因是小规模测试中 kernel 启动、固定调度开销占比很高；规模增大后 GPU 并行度和内存带宽利用率更充分，单位元素摊销开销下降。

\section{Lab 3：Histogram}

\subsection{Naive Histogram 分析}

naive 方案中每个 block 负责一个 bin，并扫描整个输入数组。因此每个输入元素会被读取 \texttt{num\_bins} 次。

当 \texttt{num\_bins=256} 且 $N=10M$ 时，总读取次数为
\[
  256\times10,000,000=2,560,000,000.
\]
若每个输入为 4 bytes，则全局内存读取流量约为
\[
  2,560,000,000\times4=10.24\ \text{GB}.
\]

该方法低效的原因是重复扫描输入数组，造成全局内存流量按 bin 数放大，并且并行方式没有充分复用每次读取的数据。

\subsection{优化后的 Histogram Kernel}

\begin{lstlisting}
__global__ void kernel(const int* input, int* histogram, int N, int num_bins) {
    __shared__ int sdata[1024];

    int idx = threadIdx.x + blockIdx.x * blockDim.x;
    int totalthread = blockDim.x * gridDim.x;

    for (int i = threadIdx.x; i < num_bins; i += blockDim.x) {
        sdata[i] = 0;
    }
    __syncthreads();

    for (int i = idx; i < N; i += totalthread) {
        atomicAdd(&sdata[input[i]], 1);
    }
    __syncthreads();

    for (int i = threadIdx.x; i < num_bins; i += blockDim.x) {
        atomicAdd(&histogram[i], sdata[i]);
    }
}
\end{lstlisting}

shared memory 上的 \texttt{atomicAdd} 比 global memory 快，因为 LDS 延迟更低、带宽更高，并且只在 block 内竞争。合并阶段每个 block 对每个 bin 至多执行一次全局 \texttt{atomicAdd}。

grid-stride loop \texttt{for (int i = idx; i < N; i += totalthread)} 的作用是让所有线程共同覆盖完整输入，即使输入规模大于总线程数也能处理，同时保持较好的负载均衡。

\subsection{实测结果与 Speedup}

两组测试的实测时间如下：

\begin{center}
\begin{tabular}{cccc}
\toprule
Test & Serial CPU (ms) & Naive GPU (ms) & Optimized GPU (ms) \\
\midrule
Test 2 medium & 2.4074  & 5.5714   & 0.1857 \\
Test 4 large  & 23.9073 & 359.2654 & 1.7784 \\
\bottomrule
\end{tabular}
\end{center}

据此可得：
\[
  \text{Test 2: Optimized vs Serial}=\frac{2.4074}{0.1857}\approx12.96\times,
  \quad
  \text{Optimized vs Naive}=\frac{5.5714}{0.1857}\approx30.00\times.
\]
\[
  \text{Test 4: Optimized vs Serial}=\frac{23.9073}{1.7784}\approx13.44\times,
  \quad
  \text{Optimized vs Naive}=\frac{359.2654}{1.7784}\approx202.01\times.
\]

naive GPU 在这两个测试中都慢于 CPU，说明重复扫描全局内存的代价很高；优化版本通过 shared memory 局部统计和一次性合并显著减少了全局内存读取与 global atomic 压力。

\subsection{Memory Efficiency}

naive 方案读取次数为 $N\times\text{num\_bins}$，优化方案每个元素只从全局内存读取一次，读取次数为 $N$。因此流量降低比例为
\[
  \frac{N\times\text{num\_bins}}{N}=\text{num\_bins}.
\]
当 \texttt{num\_bins=256} 时，理论读取流量降低 $256\times$。

\subsection{Atomic Contention}

若所有输入值都相同，例如全为 0，则大量线程会同时更新同一个 bin，atomic contention 最严重，性能会下降。实测中 Uniform 为 0.1863 ms，Gaussian 为 0.1857 ms，而 All-same(0) 为 0.2712 ms，约为 Uniform 的
\[
  \frac{0.2712}{0.1863}\approx1.46\times.
\]

若输入值在所有 bin 中均匀分布，更新会分散到多个 bin，contention 较低，因此性能通常更好。

\section{Lab 4：Parallel Reduction}

\subsection{Tree Reduction}

对 1024 个元素、block size 为 256 的输入，每个 block 内有 256 个元素参与归约，需要
\[
  \log_2 256=8
\]
步 block 内 reduction。

block 数为
\[
  \left\lceil\frac{1024}{256}\right\rceil=4,
\]
因此 block-level reduction 后剩余 4 个 partial sums。

每一步 reduction 后需要 \texttt{\_\_syncthreads()}，因为下一步会读取上一轮其它线程写入的共享内存结果。若没有同步，可能读到尚未完成更新的数据，造成错误结果。

\subsection{Atomic 数量分析}

当 $N=1,000,000$ 时：
\[
  \left\lceil\frac{1,000,000}{64}\right\rceil=15,625,
  \qquad
  \left\lceil\frac{1,000,000}{1024}\right\rceil=977.
\]
因此 block size 为 64 时最终全局 atomic 约 15,625 次；block size 为 1024 时约 977 次。

更少的全局 atomic 操作可以降低对同一输出地址的竞争，减少序列化开销，因此可能提升性能。

\subsection{算法分析}

tree reduction 第一步后，只有一半线程继续参与下一步累加，因此有 $1/2$ 的线程空闲。

warp shuffle 不需要 \texttt{\_\_syncthreads()}，因为同一个 warp 或 wavefront 内的线程按 SIMD/SIMT 方式同步执行，shuffle 指令直接在线程 lane 之间交换寄存器值，不经过共享内存。

multi-element per thread 让每个线程先在寄存器中累加多个输入元素，减少 block 数和最终 atomic 次数，同时提高每个线程的顺序内存访问工作量，有利于提升内存带宽利用率。

\subsection{实测 Reduction 结果}

naive reduction 的测试输出为：

\begin{center}
\begin{tabular}{ccc}
\toprule
Test & N & Result \\
\midrule
Small  & 10          & 4 \\
Medium & 1,000,000   & -74993790 \\
Large  & 100,000,000 & -1984960864 \\
\bottomrule
\end{tabular}
\end{center}

在 Radeon 8060S Graphics 上，对 $N=1,000,000$、block size 256 的 benchmark 结果如下：

\begin{center}
\begin{tabular}{ccc}
\toprule
Strategy & Time (ms) & Result \\
\midrule
Naive Atomic                       & 0.0312 & -74993790 \\
Tree (Shared Memory)               & 0.0250 & -74993790 \\
Tree + Warp Shuffle                & 0.0204 & -74993790 \\
Multi-Element + Tree + Shuffle     & 0.0129 & -74993790 \\
\bottomrule
\end{tabular}
\end{center}

四种策略结果一致，验证通过。相对 naive atomic 的 speedup 为：
\[
  \text{Tree}=1.25\times,\quad
  \text{Warp Shuffle}=1.53\times,\quad
  \text{Multi-Element}=2.41\times.
\]
这说明 reduction 的性能不只取决于 occupancy。减少 global atomic contention、减少同步开销、让每个线程处理多个元素以改善带宽利用，都会直接提升性能。

\section{Lab 5：Monte Carlo Integration}

\subsection{基础 Kernel 与复杂度}

\begin{lstlisting}
__global__ void montecarlo(const double* y_samples, double* result,
                           double a, double b, int n_samples) {
    int idx = blockDim.x * blockIdx.x + threadIdx.x;

    if (idx >= n_samples) return;

    double tmp = (b - a) * y_samples[idx] / n_samples;
    atomicAdd(result, tmp);
}
\end{lstlisting}

每个线程中除以 \texttt{n\_samples}，可以让每个线程直接产生最终积分中的一项贡献，最后只需要原子加到结果中。等价地，也可以先求和再统一乘以 $(b-a)/n$；后者通常与 shared-memory reduction 配合更高效。

atomic reduction 的总工作量为 $O(N)$，但由于所有线程竞争同一个全局地址，atomic 操作会产生严重序列化。

若改为 shared-memory reduction，可以让每个 block 先在共享内存中归约本 block 的局部和，再由每个 block 的一个线程对全局结果执行一次 \texttt{atomicAdd}，从而把全局 atomic 数量从 $N$ 降到 block 数。

\subsection{手工计算}

Test 1 中 $a=0,b=2,n=8$，样本为
\[
4.1805,1.7864,4.7554,3.9983,2.2786,1.1873,4.1222,1.442.
\]
样本和为
\[
  4.1805+1.7864+4.7554+3.9983+2.2786+1.1873+4.1222+1.442=23.7507.
\]
期望积分结果为
\[
  \frac{2}{8}\times23.7507=5.937675.
\]
程序输出为 5.937675，与手工计算一致。

Test 2 为 $a=b=1$、$n=1$，因此积分区间长度为 0，程序输出 0.000000 是正确的。

\subsection{Scalability 实测分析}

Monte Carlo 积分的端到端运行时间如下：

\begin{center}
\begin{tabular}{ccc}
\toprule
Test & Sample Count & Real Time (s) \\
\midrule
4 & 100,000     & 0.104 \\
7 & 10,000,000  & 0.968 \\
8 & 100,000,000 & 8.284 \\
\bottomrule
\end{tabular}
\end{center}

从 Test 4 到 Test 7，样本数量增加 $100\times$，运行时间从 0.104 s 增加到 0.968 s，约为
\[
  \frac{0.968}{0.104}\approx9.31\times.
\]
从 Test 7 到 Test 8，样本数量增加 $10\times$，运行时间约增加
\[
  \frac{8.284}{0.968}\approx8.56\times.
\]
因此端到端时间不是严格线性。小规模测试受程序启动、文件读取和固定开销影响更大；规模变大后，输入读取、host-device 传输、kernel 计算以及 atomic 累加共同决定总时间。

\subsection{Atomic Bottleneck}

若计算部分需要 1 cycle，atomic 需要 100 cycles，则每个线程约需 101 cycles，其中有效计算占比为
\[
  \frac{1}{1+100}\approx0.99\%.
\]
对于 $N=1M$，总 cycles 近似为 $1,000,000\times101$。

shared-memory reduction 可以把大部分累加移动到 block 内部，只保留每个 block 一次全局 atomic，从而显著降低全局竞争。

当 block size 为 256 且 $N=1M$ 时，全局 atomic 数为
\[
  \left\lceil\frac{1,000,000}{256}\right\rceil=3907.
\]

\subsection{数值精度}

double precision 通常提供约 15 到 16 位十进制有效数字。若改用 float，只有约 6 到 7 位有效数字。对于 $N=100M$ 的大规模累加，float 更容易出现舍入误差、低位贡献丢失以及不同累加顺序导致的结果差异。

\subsection{优化版分析}

若优化版本固定使用 256 个 blocks，则 $N=1M$ 时只需要 256 次全局 atomic。

grid-stride loop 的作用是让有限数量的线程处理任意规模输入，每个线程以固定 stride 处理多个样本，从而保持负载均衡并减少过多 block 带来的开销。

将 $(b-a)/n$ 的乘法放到 thread 0 统一完成，可以避免每个样本重复执行相同缩放操作，并让 block 内 reduction 先只处理原始样本和，减少冗余计算。

\section{Lab 6：K-Means Clustering}

\subsection{Assignment Kernel}

assignment kernel 中，每个线程负责一个点，并遍历全部 $k$ 个 centroid，计算最近中心：

\begin{lstlisting}
__global__ void find_nearest_centroids_kernel(
    const float* data_x, const float* data_y, int* labels,
    const float* centroids_x, const float* centroids_y,
    int sample_size, int k) {
    int point_id = blockIdx.x * blockDim.x + threadIdx.x;

    if (point_id < sample_size) {
        int nearest_centroid_index = 0;
        float min_distance_sq = INFINITY;

        for (int centroid_index = 0; centroid_index < k; ++centroid_index) {
            float x = centroids_x[centroid_index] - data_x[point_id];
            float y = centroids_y[centroid_index] - data_y[point_id];
            float dist = x * x + y * y;

            if (dist < min_distance_sq) {
                min_distance_sq = dist;
                nearest_centroid_index = centroid_index;
            }
        }
        labels[point_id] = nearest_centroid_index;
    }
}
\end{lstlisting}

当 $k=100$ 时，每个线程需要进行 100 次距离计算，因此单线程复杂度为 $O(k)=O(100)$，整体 assignment 阶段复杂度为 $O(Nk)$。

\subsection{Centroid Recalculation 分析}

recalculation kernel 先使用 shared memory atomics，再合并到 global memory，是为了把大量点的累加先限制在 block 内，降低 global atomic contention。每个 block 最后只需要按 centroid 数执行合并，减少对全局内存的竞争。

若某个 centroid 没有任何点分配给它，代码会保留上一轮 centroid 坐标，避免除以 0，并保证空 cluster 不会产生无效值。

\subsection{Compute Intensity}

对于 $N=50,000$、$k=50$：

\begin{itemize}
  \item 每个点的距离计算次数：50。
  \item 总距离计算次数：$50,000\times50=2,500,000$。
  \item update 阶段 shared memory atomics：每个点对 $x$、$y$ 和 count 各一次，共 $50,000\times3=150,000$ 次。
  \item 若 block size 为 256，block 数为 $\lceil 50,000/256\rceil=196$。每个 block 对每个 cluster 执行 $x$、$y$ 和 count 三类 global atomic，因此 global atomics 为 $196\times50\times3=29,400$ 次。
\end{itemize}

Test 4 的实测配置为 $N=50,000$、$k=50$、max iterations 为 100，端到端运行时间为 0.113 s。每轮迭代的主导计算通常是 assignment 阶段的 $N\times k$ 次距离计算，尤其当 $k$ 较大时更明显。本测试中每轮 assignment 需要 2,500,000 次距离计算，因此它是主要计算来源；update 阶段则主要受 atomic 累加和全局合并影响。

\subsection{Memory Access Optimization}

centroid 会被所有线程反复读取，可以将 centroid 分批加载到 shared memory 或利用 constant/cache-friendly 访问模式，减少重复 global memory 读取。若 $k$ 较小，也可以提高 centroid 数据的缓存命中率。

当 $k=1000$ 时，centroid 数组变大，单个线程需要遍历更多中心，缓存工作集增大，cache miss 和内存带宽压力都会上升，同时 assignment 阶段计算量也增加到 $O(1000)$ 每点。

\subsection{Convergence}

本实现的收敛阈值为
\[
  dx^2+dy^2<1.0\times10^{-8}.
\]
当全部 $k$ 个 centroid 都满足该条件时提前终止。

检查收敛比固定运行 \texttt{max\_iterations} 更重要，因为许多数据集会提前稳定。提前停止可以减少无意义迭代，节省时间和能耗。

double-buffering 通过交换上一轮 centroid 和新 centroid 的指针，避免在同一数组中边读边写造成数据覆盖，同时减少额外拷贝开销。

\section{总结}

六个实验的核心思想分别覆盖了一维索引、二维索引与矩阵转置、直方图 atomic 优化、并行归约、Monte Carlo 积分归约以及 K-Means 迭代聚类。总体规律是：GPU 程序不仅要保证线程覆盖完整数据，还需要关注 coalesced memory access、shared memory 局部归约、atomic contention、wavefront 对齐以及 occupancy 与实际瓶颈之间的关系。

\end{document}
